\documentclass[12pt,a4paper]{article}

% --- Настройка языка и шрифтов для LuaLaTeX ---
\usepackage{fontspec}
\defaultfontfeatures{Ligatures=TeX,Scale=MatchLowercase}

\setmainfont{Alice-Regular}[
    Path=font/,
    Extension=.ttf,
    BoldFont = Alice-Regular,
    ItalicFont = Alice-Regular,
    BoldItalicFont = Alice-Regular,
    BoldFeatures = {FakeBold=1.6},
    ItalicFeatures = {FakeSlant=0.2},
    BoldItalicFeatures = {FakeBold=1.6,FakeSlant=0.2}
]

\newfontfamily\cyrillicfont{Alice-Regular}[
    Path=font/,
    Extension=.ttf,
    BoldFont = Alice-Regular,
    ItalicFont = Alice-Regular,
    BoldItalicFont = Alice-Regular,
    BoldFeatures = {FakeBold=1.6},
    ItalicFeatures = {FakeSlant=0.2},
    BoldItalicFeatures = {FakeBold=1.6,FakeSlant=0.2}
]

\usepackage{polyglossia}
\setmainlanguage{russian}
\setsansfont{TeX Gyre Heros}
\setmonofont{TeX Gyre Cursor}
\AtBeginDocument{\renewcommand{\textnumero}{{\sffamily №}}}

\usepackage{graphicx}
\usepackage{amsmath, amsfonts, amssymb, amsthm}
\usepackage{mathtools}
\usepackage[protrusion=false]{microtype}
\usepackage{unicode-math}
\usepackage{array}
\usepackage{float}
\usepackage{caption}
\usepackage{enumitem}

\usepackage[paper=letterpaper,
            top=20mm,
            bottom=20mm,
            left=10mm,
            right=10mm,
            footskip=0.4in 
]{geometry}

% --- Колонтитулы ---
\usepackage{fancyhdr}
\pagestyle{fancy}
\fancyhf{}
\setlength{\headheight}{15pt}

\renewcommand{\headrulewidth}{0.4pt}
\renewcommand{\footrulewidth}{0.4pt}

\newcommand{\DocTitle}{Методы борьбы с ИИ}
\newcommand{\Discipline}{Социология}

\begin{document}

% Движение лудитов (движение против отказа от станков на тяге паровых машин)
% 1. Контекст и проблематика ограничительных мер ИИ
% Необходимы реформы в ваккуме института образования
% 2. Институциональные ограничения ИИ (на уровне института образования)
% 3. На уровне конкретных организаций (например, университетов)

\fancyhead[L]{\DocTitle}
\fancyhead[R]{\Discipline}
\fancyfoot[C]{\thepage}

\section{Введение и проблематика}

Борьба с ИИ в образовательном процессе является во многом сложным дискуссионным 
вопросом, в котором есть множество
решений, которые указывают на необходимость развала архаичных институтов образования
в пользу децентрализованного формата с упрощённой сертификацией знаний, что тем не 
менее не является реалистичным решением для социальных стран, коими являются
все страны Европы.

Поэтому в целях обсуждения реальных путей устранения проблемы злоупотребления студентами
ИИ в процессе получения образования мы рассмотрим решения, которые будут направлены на 
сохранение существующих институтов образования или по крайней мере гарантируют 
правопреемственность существующих институтов образования.

Также хотелось бы дополнительно разграничить борьбу с использованием ИИ в образовательном 
процессе и борьбу со списыванием, которая обладает несколько иной проблематикой 
и требует отдельного рассмотрения, поэтому в рамках данной работы не будет 
рассматриваться проблема использования ИИ на конвенциональных контрольных мероприятиях
или будет упоминаться без детального рассмотрения.

\section{Методика борьбы через санкции на использование}

Одним из первых методов борьбы с использованием ИИ, который имеет некоторую 
практику реализации в высших учебных заведениях, является борьба с использованием ИИ 
через санкции, например, занижение баллов за текст, написанный с помощью ИИ. 

Упор в тексты связан с проблемой обнаружения использования ИИ в иных форматах работ, 
при этом в то же время любой формат работы легко трансформируется в текстовую форму.  
Также различные форматы, например, тестовые, обладают высоким потенциалом для внутреннего 
контроля учебным заведением на использование ИИ посредством процедуры её проведения, 
что во многом недоступно текстовым работам из-за времени выполнения, кое слишком велико
для объёма учебных часов, поэтому неизбежно уходит на выполнение вне стен учебного 
заведения и тем более вне контроля преподавателя.

\subsection{Практическая реализация метода}

Рассмотрим опыт ведущих университетов США, включая, например, университет Вандербильта, в
которых был реализован ранее описанный метод борьбы с ИИ. 

На практике проверка текста
на использование ИИ осуществлялась через специализированное ПО, которое посредством
анализа перплексии и бёрстности текстового материала, оценивая степень вероятности и 
вариативности следующих друг за другом токенов, давало оценку использования ИИ в тексте. 
Однако этот инструментарий быстро продемонстрировал концептуальную несостоятельность. 
Алгоритмы оказались неспособны надежно отличать структурированный стиль человека от 
генерации языковой модели, что породило проблему систематических ложноположительных 
срабатываний. 

Масштабное исследование Стэнфордского университета зафиксировало, что 
ведущие детекторы классифицировали реальные студенческие эссе TOEFL, написанные 
авторами, для которых английский не является родным языком, как сгенерированные машиной 
в среднем в $61{,}2\%$ случаев, причём подавляющее большинство работ было помечено хотя 
бы одной программой. Одновременно с этим студенты начали массово обходить алгоритмы с 
помощью ручного перефразирования, изменения пунктуации и использования промежуточных 
сервисов синтаксической модификации. 

Высокие риски несправедливых обвинений и 
отсутствие юридической силы показателей заставили официально отказаться от 
использования детекторов как основания для дисциплинарных мер. Иные метрики оказались
субъективнее и фальсифицируемее, что не позволило пересмотреть неудачную практику 
реализации метода борьбы через санкции.

\section{Метод институционального запрета}

Когда не работают ограничения, в ход идут запретительные меры. Метод институционального 
запрета предполагает запрет на использование ИИ в образовательном процессе на уровне
института образования. 

Целеполагание метода заключается в том, чтобы на уровне 
учебного заведения создать условия, при которых использование ИИ в образовательном
процессе будет невозможно, притом для обеспечения работоспособности метода необходимо 
создать условия, которые принудят студентов к обучению исключительно в стенах учебного 
заведения.

Обучение на дому и выполнение в домашних условиях любого рода учебных работ 
сведёт инфраструктуру и механизм метода к формальным мерам, которые не возымеют 
желаемого эффекта. А именно, из целеполагания метода следует, что студент всегда 
предпочтёт использовать ИИ при его наличии, а вне стен учебного заведения он будет 
в наличии, что создаст условия, в которых всё обучение проходит с использованием ИИ, 
а успеваемость на контрольных мероприятиях будет проверяться вне привычных условий, что 
заметно понизит успеваемость, как и создаст эффект точечной подготовки, при котором
обучение заточено под задачи экзамена и не несёт утилитарного смысла, заложенного 
учебной программой. 

\subsection{Практическая реализация метода}

Рассмотрим опыт школьных округов Нью-Йорка и Лос-Анджелеса — блокировка доступа к 
платформам ИИ на территории учебных заведений. Учащиеся без труда переходили на 
персональные устройства, мобильный интернет и сторонние прокси-серверы. Статистические 
замеры подтвердили масштаб проблемы: свыше $85\%$ студентов применяли генеративные 
системы в учёбе, а около четверти полностью заменяли ими выполнение письменных 
домашних заданий. Метод в условиях защиты принципов 
частной собственности и приватности личной жизни оказался неэффективным или, скорее, 
нереализуемым на практике в полной мере. 

\subsection{Исключительные случаи}

Хотя в случаях, когда дисциплина позволяет 
отойти от благ цифровой инфраструктуры, некоторые учебные заведения успешнее реализуют
метод институционального запрета посредством ренессанса аналоговых рукописных экзаменов 
в традиционных синих тетрадях под строгим очным надзором преподавателей. 

Ведущие центры, включая Принстон и Стэнфорд, вернулись к обязательному очному 
прокторингу, рассматривая изоляцию учащихся от сетевых устройств как единственный 
действенный способ верификации остаточных знаний. Параллельно институты внедрили 
концепцию процессуального оценивания, где предметом проверки становится не финальный 
текст, а последовательность работы: промежуточные черновики, фиксация версионности 
файлов, источниковедческие записки и очные собеседования в формате сократического 
диалога, на которых студент обязан подтвердить логику и владение материалом 
собственной работы.

\section{Метод частичного ограничения применения}

Ограничения показали свою неэффективность, запретительные меры оказались во многом 
формальностью, поэтому в качестве альтернативы был выработан новый метод, который во 
многом является компромиссом между консервативными и прогрессивными подходами. 
Метод частичного ограничения применения предполагает использование ИИ в 
образовательном процессе, но регламентирует область его применения, исключая 
чувствительные к образовательному процессу моменты. 

Целеполагание метода предполагает, 
что студент при наличии процедуры, позволяющей легальное применение ИИ в автоматизации 
рутинных задач, которые не несут ключевой для дисциплины образовательной ценности, 
будет склонен не обесценивать образовательный процесс и применять ИИ исключительно 
в рамках своего удобства, увеличивая свой комфорт и эффективность обучения.

\subsection{Практическая реализация метода}

Школьные округа Нью-Йорка и Лос-Анджелеса, которые ранее пытались реализовать метод 
институционального запрета, в итоге перешли к частичному ограничению применения ИИ.
В целом подход является доминирующей парадигмой в образовательных институтах, так 
как не требует радикальных изменений в инфраструктуре и хорошо согласуется с 
гражданскими свободами.

Но у 
метода есть и подводные камни — устранение когнитивного сопротивления из процесса 
обучения разрушает биологические предпосылки формирования глубокого интеллекта и 
систематической эрудиции, в свою очередь когнитивное сопротивление является 
ключевым фактором дискомфорта образовательного
процесса. Если учитывать, что обратный эффект Флинна\footnote{Эффект Флинна — это 
явление прироста метрики естественного интеллекта IQ с каждым новым поколением. 
На протяжении большей части XX столетия психометрические замеры фиксировали 
стабильный межпоколенческий прирост показателей коэффициента интеллекта (IQ). 
Однако лонгитюдные когортные исследования призывников в Дании, Норвегии, 
Великобритании и Финляндии выявили стагнацию и выраженный спад показателей IQ у 
поколений, родившихся после середины 1990-х годов. Данные международных замеров 
PISA зафиксировали беспрецедентный синхронный спад функциональной грамотности по 
математике, чтению и естественным наукам по всем странам Организации 
экономического сотрудничества и развития. Выступая перед комитетом Сената США, 
нейробиолог доктор Джаред Куни Хорват представил доказательства того, что снижение 
успеваемости напрямую коррелирует с распространением планшетов, ноутбуков и 
фрагментированного цифрового контента в учебных аудиториях, подорвавшего 
концентрацию внимания и спровоцировавшего феномен функциональной неграмотности.} 
фактически реализуется, то подход оказался неэффективным на практике и содержит 
ошибку в целеполагании — привнесение комфорта в процесс обучения не является 
эффективной стратегией или по крайней мере не является абсолютной.

\section{Заключение}

Каждый рассмотренный метод имеет ключевые противоречия между гражданскими свободами 
и образовательной эффективностью. В связи с этим реализация базового права на 
образование в условиях широкого распространения ИИ включает в себя выбор между 
сохранением эффективности образования и обеспечением комфорта его получения.

\end{document}
